\section{Respuesta social: qué hacer después de la deflación del trabajo de oficina}

La parte más sombría de este episodio es la deflación del trabajo de oficina y el periodo de desempleo que abre. En lo técnico, los modelos de Coding acortan los ciclos de desarrollo; en lo organizacional, unas pocas personas pueden completar más tareas; en lo económico, el precio del trabajo de oficina se volverá a fijar. La pregunta no es «¿reemplazará la IA a las personas?», sino si la sociedad puede ofrecer un nuevo lugar a quienes sean reemplazados o reorganizados.

\subsection{Tres grupos que recibirán el impacto primero}

\noindent\begin{tabularx}{\textwidth}{p{0.22\textwidth}p{0.34\textwidth}X}
\toprule
Grupo & Por qué lo afecta primero & Dirección de respuesta \\
\midrule
Trabajadores del conocimiento con tareas repetitivas & Su trabajo puede expresarse en texto, tablas y código & Pasar a la definición de tareas, la revisión, la comprensión del cliente y el diseño de procesos. \\
Programadores principiantes & El modelo puede generar gran parte del código de implementación & Aprender arquitectura, pruebas, review y comprensión de sistemas. \\
Coordinadores de mando medio & Las herramientas pueden automatizar el reenvío de información y el seguimiento del avance & Pasar al juicio, la asignación de recursos y la resolución de conflictos entre equipos. \\
\bottomrule
\end{tabularx}

\begin{knowledgebox}{El verdadero objetivo de la recapacitación}
La recapacitación no debe limitarse a enseñar a «usar herramientas de IA»; debe formar a las personas para definir problemas, verificar resultados, organizar flujos de trabajo y comprender el negocio. Las herramientas cambiarán; el criterio y la estructura de responsabilidades cambian más despacio.
\end{knowledgebox}

\subsection{Resumen de la sección}

La deflación del trabajo de oficina resulta de que la tecnología se difunde más rápido de lo que las organizaciones pueden absorberla. La respuesta social no puede depender solo del aprendizaje individual: también hace falta que la organización de las empresas y el diseño institucional redistribuyan los beneficios del aumento de productividad.
